\section{Algorithms}

\subsection{FDCheck: Adding a File}

When a file is added, its size is computed (\texttt{file\_size}) and used as the key to store the file in a \textbf{Hashmap}.
If another file with the same size is added, a new \textbf{Hashmap} is created that indexes files by their content.
For each file, a hash is computed (\texttt{file\_hash\_by\_path\_ptr}) and used for fast lookup in the table;
in case of hash collisions, a byte-by-byte comparison of the files is performed (\texttt{file\_equal\_by\_path\_ptr}).
The value in this map is a \textbf{Hashset} that stores file paths,
also using a hash (\texttt{string\_hash\_by\_ptr}) derived from the string
and performing string comparison (\texttt{string\_equal\_by\_ptr}) when necessary.


\subsection{Tabulation Hashing}

Files and strings are hashed using tabulation hashing with a table of $32 \times 256$ random 64-bit (\texttt{size\_t}) values.
The table is initialized once on the first call to the hash function.

\subsection{FDDump: Retrieving Results}

To obtain the results, we iterate over all hashmaps and compute the total number of entries.
Then we allocate a table of the required size and iterate again to populate it.
